\documentclass[11pt]{article}

\usepackage[a4paper,margin=1in]{geometry}
\usepackage{amsmath,amssymb,amsthm,mathtools}
\usepackage{enumitem}
\usepackage{microtype}
\usepackage{tcolorbox}

\setlength{\parskip}{0.5em}
\setlength{\parindent}{0pt}

\newcommand{\E}{\mathbb{E}}
\newcommand{\PP}{\mathbb{P}}
\newcommand{\R}{\mathbb{R}}
\newcommand{\F}{\mathcal{F}}
\newcommand{\G}{\mathcal{G}}
\newcommand{\1}{\mathbf{1}}
\newcommand{\vol}{\mathrm{vol}}

\theoremstyle{plain}
\newtheorem{claim}{Claim}

\title{IE 622 \\ Detailed Solutions --- Practice Problem Set 12}
\author{}
\date{}

\begin{document}
\maketitle

% ======================================================================
\section*{Key Concepts Used in This Problem Set}

This section summarises the central definitions and theorems that appear
throughout the problems. Reading this first will make every proof below
easier to follow.

% ------ Concept 1 -------------------------------------------------------
\subsection*{1. Conditional Expectation}

\paragraph{Definition.}
Let $(\Omega,\F,P)$ be a probability space, $X$ an integrable r.v., and
$\G\subseteq\F$ a sub-$\sigma$-algebra. The \emph{conditional expectation}
$E[X\mid\G]$ is the (a.s.\ unique) random variable $Z$ such that:
\begin{enumerate}
  \item $Z$ is $\G$-measurable.
  \item $\displaystyle\int_G Z\,dP = \int_G X\,dP$ for every $G\in\G$
        \quad(the \emph{partial-averaging} property).
\end{enumerate}

\paragraph{Key properties used in the solutions.}
\begin{itemize}
  \item \textbf{Tower property}: If $\mathcal{A}\subseteq\mathcal{B}\subseteq\F$,
        then $E\bigl[E[X\mid\mathcal{B}]\mid\mathcal{A}\bigr]=E[X\mid\mathcal{A}]$ a.s.
  \item \textbf{Pulling out knowns}: If $Y$ is $\G$-measurable and $E|XY|<\infty$,
        then $E[XY\mid\G]=Y\cdot E[X\mid\G]$ a.s.
  \item \textbf{Measurability}: If $X$ is $\G$-measurable, then $E[X\mid\G]=X$ a.s.
  \item \textbf{Independence}: If $X\perp\G$, then $E[X\mid\G]=E[X]$ a.s.
\end{itemize}

% ------ Concept 2 -------------------------------------------------------
\subsection*{2. Filtrations}

A \emph{filtration} is an increasing sequence of $\sigma$-algebras
$\F_1\subseteq\F_2\subseteq\cdots\subseteq\F$. Intuitively, $\F_n$ represents
the information available up to time $n$. The \emph{natural filtration} of a
process $\{X_n\}$ is $\F_n^X=\sigma(X_1,\ldots,X_n)$.

% ------ Concept 3 -------------------------------------------------------
\subsection*{3. Martingales, Supermartingales, and Submartingales}

\paragraph{Definition.}
An adapted, integrable process $\{M_n,\F_n\}$ is a:
\begin{itemize}
  \item \textbf{Martingale} if $E[M_{n+1}\mid\F_n]=M_n$ a.s.\ for all $n$.
  \item \textbf{Supermartingale} if $E[M_{n+1}\mid\F_n]\le M_n$ a.s.\ for all $n$.
  \item \textbf{Submartingale} if $E[M_{n+1}\mid\F_n]\ge M_n$ a.s.\ for all $n$.
\end{itemize}

\paragraph{Intuition.}
A martingale models a \emph{fair game}: your expected future wealth equals your
current wealth. A supermartingale models an \emph{unfair game} (losing on average),
and a submartingale a winning one. The classic example: $M_n=$ fortune after $n$
steps of a fair coin-flip game is a martingale.

\paragraph{Standard martingales used in this set.}
\begin{enumerate}
  \item $\{X_n\}$ = symmetric random walk ($E[X_{n+1}\mid\F_n]=X_n$).
  \item $\{X_n^2-n\}$ = squared symmetric random walk minus time (since
        $\mathrm{Var}(Z_i)=1$).
  \item $\{\mu^{-n}X_n\}$ in Galton--Watson process.
  \item $\{R_n/(n+r+b)\}$ in P\'olya's urn.
\end{enumerate}

% ------ Concept 4 -------------------------------------------------------
\subsection*{4. The Optional Stopping Theorem (OST)}

\paragraph{Statement.}
Let $\{M_n,\F_n\}$ be a martingale and $T$ a stopping time. Under any of the
following sufficient conditions:
\begin{enumerate}[label=(\alph*)]
  \item $T$ is bounded ($T\le N$ for some fixed $N$), or
  \item The stopped process $\{M_{n\wedge T}\}$ is dominated by an integrable
        r.v.\ and $T<\infty$ a.s., or
  \item $E[T]<\infty$ and increments $|M_{n+1}-M_n|$ are bounded,
\end{enumerate}
we have $E[M_T]=E[M_0]$.

\paragraph{Why it's powerful.}
OST lets us evaluate expectations at random (stopping) times by relating them to
the initial value. In Problems 8(a) and 8(b), this is the core tool.

% ------ Concept 5 -------------------------------------------------------
\subsection*{5. Markov Property of DTMCs}

A DTMC $\{X_n\}$ satisfies the \emph{Markov property}:
\[
  P(X_{n+1}=y\mid X_0,\ldots,X_n) = P(X_{n+1}=y\mid X_n) = p_{X_n,y}.
\]
Equivalently, for any measurable $g$,
\[
  E[g(X_{n+1})\mid\mathcal{F}_n] = E[g(X_{n+1})\mid X_n]
  = \sum_{y}p_{X_n,y}\,g(y).
\]
This is used in Problems 3, 5, and 6 to simplify conditional expectations.

% ------ Concept 6 -------------------------------------------------------
\subsection*{6. Generator of a Markov Chain}

For a DTMC with transition matrix $P$, the \emph{(discrete) generator} (or
Dynkin operator) acting on $f:\mathcal{S}\to\R$ is
\[
  (\mathcal{L}f)(x) = \sum_{y}p_{x,y}(f(y)-f(x)) = E[f(X_1)-f(X_0)\mid X_0=x].
\]
This measures the \emph{expected drift} of $f$ in one step. When $(\mathcal{L}f)(x)=0$
for all $x$, $f$ is called \emph{harmonic}, and $\{f(X_n)\}$ is a martingale.

% ------ Concept 7 -------------------------------------------------------
\subsection*{7. Superharmonic Functions and the Mean Value Property}

A function $f:\R^d\to\R$ is \emph{superharmonic} if it satisfies the
sub-mean-value property:
\[
  f(x) \ge \frac{1}{\vol(B(x,r))}\int_{B(x,r)}f(y)\,dy \quad\forall\,x,r>0.
\]
Intuitively: the value of $f$ at any point is at least the average of $f$ over
any surrounding ball. Harmonic functions (like $f(x)=x$ in 1D, or $f(x)=\log|x|$
in 2D for $x\neq0$) satisfy equality. The connection to supermartingales is
explored in Problem 5.

% ------ Concept 8 -------------------------------------------------------
\subsection*{8. Doob Decomposition (background)}

Any adapted integrable process $\{X_n\}$ can be uniquely written as
$X_n=M_n+A_n$ where $\{M_n\}$ is a martingale and $\{A_n\}$ is a predictable
process (i.e., $A_n$ is $\F_{n-1}$-measurable) with $A_0=0$. In Problem 3, the
sum $\sum_{i=0}^{n-1}(\mathcal{L}f)(X_i)$ serves as the predictable compensator
$A_n$, and $Y_n=f(X_n)-f(X_0)-A_n$ is the martingale part.

% ======================================================================
\section*{Problem 1}

\paragraph{Problem statement.}
Let $(\Omega,\F,P)$ be a probability space. Suppose $\{B_i\}_{i\ge1}\subset\F$ is
a disjoint partition of $\Omega$ (i.e., $\Omega=\bigcup_{i=1}^{\infty}B_i$ and
$B_i\cap B_j=\emptyset$ for $i\neq j$) with $P(B_i)>0$ for all $i$.
Let $\G=\sigma(B_1,B_2,\ldots)$. For an $\F$-measurable, integrable random
variable $X$, show that
\[
  E[X\mid\G] = \frac{E[X\mathbf{1}_{B_i}]}{P(B_i)}
  \quad\text{on }B_i.
\]

\paragraph{Background: definition of conditional expectation.}
Recall that $Z=E[X\mid\G]$ is the \emph{unique} (a.s.) $\G$-measurable,
integrable random variable satisfying
\[
  \int_G Z\,dP = \int_G X\,dP \quad\text{for every }G\in\G.
\]
We will verify that $Z:=\sum_{i=1}^{\infty}\frac{E[X\mathbf{1}_{B_i}]}{P(B_i)}\mathbf{1}_{B_i}$
satisfies both conditions.

\paragraph{Solution.}

\textbf{Step 1 — Describing the $\sigma$-algebra $\G$.}
Because $\{B_i\}$ is a partition of $\Omega$, the atoms of $\G=\sigma(B_1,B_2,\ldots)$
are exactly the sets $B_i$ themselves. Any $G\in\G$ is a (possibly empty, possibly
infinite) union $G=\bigcup_{j\in J}B_j$ for some $J\subseteq\mathbb{N}$. This
follows because: (a) every such union is a member of $\G$; (b) the class of such
unions is closed under complements ($\Omega\setminus\bigcup_{j\in J}B_j
=\bigcup_{j\notin J}B_j$) and countable unions, hence is itself a $\sigma$-algebra
containing every $B_i$, making it equal to $\G$.

\textbf{Step 2 — Defining the candidate $Z$.}
For each $\omega\in\Omega$ there is a unique index $i(\omega)$ with
$\omega\in B_{i(\omega)}$. Define
\[
  Z(\omega) := \frac{E[X\mathbf{1}_{B_{i(\omega)}}]}{P(B_{i(\omega)})},
\]
or equivalently,
\[
  Z = \sum_{i=1}^{\infty} c_i\,\mathbf{1}_{B_i}, \quad
  c_i := \frac{E[X\mathbf{1}_{B_i}]}{P(B_i)}.
\]
Note: $E[X\mathbf{1}_{B_i}]$ is a real number (it is a finite constant because
$X$ is integrable and $\mathbf{1}_{B_i}\le1$, so
$|E[X\mathbf{1}_{B_i}]|\le E[|X|\mathbf{1}_{B_i}]<\infty$). Integrability of $Z$:
\[
  E|Z|
  = \sum_{i=1}^{\infty}|c_i|\,P(B_i)
  = \sum_{i=1}^{\infty}|E[X\mathbf{1}_{B_i}]|
  \le \sum_{i=1}^{\infty}E[|X|\mathbf{1}_{B_i}]
  = E|X| < \infty.
\]

\textbf{Step 3 — Measurability.}
Each indicator $\mathbf{1}_{B_i}$ is $\G$-measurable (since $B_i\in\G$).
A countable linear combination of $\G$-measurable functions (with deterministic
coefficients $c_i$) is $\G$-measurable. Hence $Z$ is $\G$-measurable.

\textbf{Step 4 — Partial averaging (the key verification).}
Take any $G\in\G$. By Step 1, $G=\bigcup_{j\in J}B_j$ for some $J\subseteq\mathbb{N}$.
The sets $\{B_j\}_{j\in J}$ are disjoint, so
\begin{align*}
  \int_G Z\,dP
  &= E[Z\,\mathbf{1}_G]
   = E\!\left[\left(\sum_{i=1}^{\infty}c_i\,\mathbf{1}_{B_i}\right)
              \mathbf{1}_{\bigcup_{j\in J}B_j}\right].
\end{align*}
Since $\mathbf{1}_{B_i}\cdot\mathbf{1}_{\bigcup_{j\in J}B_j}=\mathbf{1}_{B_i}$
if $i\in J$ and $0$ otherwise, the sum collapses to $j\in J$ only:
\[
  = E\!\left[\sum_{j\in J}c_j\,\mathbf{1}_{B_j}\right]
  = \sum_{j\in J}c_j\,E[\mathbf{1}_{B_j}]
  \quad\text{(Fubini/MCT, since terms are non-negative after splitting $X=X^+-X^-$)}
\]
\[
  = \sum_{j\in J}\frac{E[X\mathbf{1}_{B_j}]}{P(B_j)}\cdot P(B_j)
  = \sum_{j\in J}E[X\mathbf{1}_{B_j}]
  = E\!\left[X\sum_{j\in J}\mathbf{1}_{B_j}\right]
  = E[X\,\mathbf{1}_G]
  = \int_G X\,dP.
\]
(The interchange of expectation and sum is justified because
$\sum_{j\in J}E[|X|\mathbf{1}_{B_j}]\le E|X|<\infty$.)

\textbf{Conclusion.}
$Z$ is $\G$-measurable, integrable, and satisfies the partial-averaging property
for every $G\in\G$. By uniqueness of conditional expectation (a.s.),
\[
  E[X\mid\G] = Z = \frac{E[X\mathbf{1}_{B_i}]}{P(B_i)}\quad\text{on }B_i.
  \qquad\square
\]

% ======================================================================
\section*{Problem 2}

\paragraph{Problem statement.}
Suppose $\{X_n\}$ is a martingale with respect to $\{\F_n\}$. Let
$\G_n=\sigma(X_1,\ldots,X_n)$. Show that $\{X_n\}$ is a martingale with
respect to $\{\G_n\}$.

\paragraph{Key observations.}
\begin{enumerate}
  \item $\G_n\subseteq\F_n$: each $X_k$ for $k\le n$ is $\F_k\subseteq\F_n$-measurable,
        so $\sigma(X_1,\ldots,X_n)\subseteq\F_n$.
  \item The filtration $\{\G_n\}$ is increasing: $\G_n\subseteq\G_{n+1}$ since
        $\sigma(X_1,\ldots,X_n)\subseteq\sigma(X_1,\ldots,X_n,X_{n+1})$.
\end{enumerate}

\paragraph{Solution.}

\textbf{(i) Integrability.} Since $\{X_n\}$ is a martingale w.r.t.\ $\{\F_n\}$,
$E|X_n|<\infty$ for all $n$. So $X_n$ is integrable.

\textbf{(ii) Adaptedness.} By definition of $\G_n=\sigma(X_1,\ldots,X_n)$, the
random variable $X_n$ is $\G_n$-measurable (it is one of the generating variables).

\textbf{(iii) Martingale property.}
We must show $E[X_{n+1}\mid\G_n]=X_n$ a.s.

\underline{Step 1}: Apply the tower property. Since $\G_n\subseteq\F_n$,
\[
  E[X_{n+1}\mid\G_n]
  = E\!\bigl[E[X_{n+1}\mid\F_n]\mid\G_n\bigr].
\]
(The tower property states: if $\mathcal{A}\subseteq\mathcal{B}$ are
sub-$\sigma$-algebras, then $E[E[Y\mid\mathcal{B}]\mid\mathcal{A}]
=E[Y\mid\mathcal{A}]$.)

\underline{Step 2}: Since $\{X_n\}$ is a martingale w.r.t.\ $\{\F_n\}$,
\[
  E[X_{n+1}\mid\F_n] = X_n\quad\text{a.s.}
\]

\underline{Step 3}: Substituting,
\[
  E[X_{n+1}\mid\G_n]
  = E[X_n\mid\G_n].
\]

\underline{Step 4}: $X_n$ is $\G_n$-measurable (Step (ii)), so
$E[X_n\mid\G_n]=X_n$ a.s. (a measurable random variable is its own conditional
expectation given its own $\sigma$-algebra).

Combining:
\[
  E[X_{n+1}\mid\G_n] = X_n\quad\text{a.s.}
\]
Therefore $\{X_n\}$ is a martingale with respect to $\{\G_n\}$. $\square$

\paragraph{Remark.}
The key structural fact is that $\G_n\subseteq\F_n$: passing to a \emph{coarser}
filtration preserves the martingale property. The $\G_n$-martingale property then
follows by smoothing over the extra information in $\F_n$.

% ======================================================================
\section*{Problem 3 (Dynkin's Martingale)}

\paragraph{Problem statement.}
Let $\{X_n\}$ be a DTMC on state space $\mathcal{S}$ with transition matrix $P$.
Define the \emph{generator} operator
\[
  (\mathcal{L}f)(x) := \sum_{y\in\mathcal{S}}p_{x,y}(f(y)-f(x))
  = E[f(X_1)-f(X_0)\mid X_0=x].
\]
Set $Y_0=0$ and, for $n\ge1$,
\[
  Y_n := f(X_n) - f(X_0) - \sum_{i=0}^{n-1}(\mathcal{L}f)(X_i).
\]
Show that $\{Y_n\}$ is a martingale with respect to $\mathcal{H}_n:=\sigma(X_0,\ldots,X_n)$.

\paragraph{Intuition.}
The sum $\sum_{i=0}^{n-1}(\mathcal{L}f)(X_i)$ is the \emph{predictable compensator}
of the process $\{f(X_n)\}$, i.e., it removes the expected drift so that the
residual is a martingale. This construction mirrors the Doob decomposition.

\paragraph{Solution.}

\textbf{(i) Adaptedness.}
$Y_n$ depends only on $X_0,X_1,\ldots,X_n$, hence is $\mathcal{H}_n$-measurable.

\textbf{(ii) Integrability.}
Assume $f$ is bounded (or that $E|f(X_n)|<\infty$ for all $n$; since $(\mathcal{L}f)(x)
=\sum_y p_{x,y}f(y)-f(x)$ is a finite linear combination of values of $f$, finiteness
of $E|f(X_k)|$ implies finiteness of $E|(\mathcal{L}f)(X_k)|$). Then $E|Y_n|<\infty$.

\textbf{(iii) Martingale property.}
Compute the increment:
\[
  Y_{n+1} - Y_n
  = \Bigl[f(X_{n+1})-f(X_0)-\sum_{i=0}^{n}(\mathcal{L}f)(X_i)\Bigr]
    - \Bigl[f(X_n)-f(X_0)-\sum_{i=0}^{n-1}(\mathcal{L}f)(X_i)\Bigr]
  = f(X_{n+1}) - f(X_n) - (\mathcal{L}f)(X_n).
\]
Now take the conditional expectation given $\mathcal{H}_n$:
\begin{align*}
  E[Y_{n+1}-Y_n\mid\mathcal{H}_n]
  &= E[f(X_{n+1})\mid\mathcal{H}_n] - f(X_n) - (\mathcal{L}f)(X_n).
\end{align*}
By the Markov property of $\{X_n\}$,
\[
  E[f(X_{n+1})\mid\mathcal{H}_n]
  = E[f(X_{n+1})\mid X_n]
  = \sum_{y\in\mathcal{S}}p_{X_n,y}\,f(y).
\]
Note that
\[
  (\mathcal{L}f)(X_n)
  = \sum_{y\in\mathcal{S}}p_{X_n,y}(f(y)-f(X_n))
  = \sum_{y\in\mathcal{S}}p_{X_n,y}f(y) - f(X_n)\sum_{y\in\mathcal{S}}p_{X_n,y}
  = \sum_{y\in\mathcal{S}}p_{X_n,y}f(y) - f(X_n),
\]
since $\sum_y p_{X_n,y}=1$. Therefore,
\[
  E[f(X_{n+1})\mid\mathcal{H}_n] = (\mathcal{L}f)(X_n) + f(X_n).
\]
Substituting back:
\begin{align*}
  E[Y_{n+1}-Y_n\mid\mathcal{H}_n]
  &= \bigl[(\mathcal{L}f)(X_n)+f(X_n)\bigr] - f(X_n) - (\mathcal{L}f)(X_n)
   = 0.
\end{align*}
Therefore $E[Y_{n+1}\mid\mathcal{H}_n]=Y_n$ a.s., proving $\{Y_n\}$ is a
martingale. $\square$

% ======================================================================
\section*{Problem 4 (Galton--Watson Branching Process)}

\paragraph{Problem statement.}
Let $\{Z_i^n\}$ be i.i.d.\ nonneg.\ integer-valued r.v.'s with mean $\mu>0$.
Define $X_0=1$ and $X_{n+1}=\sum_{i=1}^{X_n}Z_i^{n+1}$ for $n\ge0$.
Let $\mathcal{F}_n=\sigma(Z^1,Z^2,\ldots,Z^n)$, where $Z^k=(Z_1^k,Z_2^k,\ldots)$.
Show that $\{M_n\}:=\{\mu^{-n}X_n, n\ge0\}$ is a martingale w.r.t.\ $\{\mathcal{F}_n\}$.

\paragraph{Setup and notation.}
$X_n$ = population size at generation $n$.
$Z_i^{n+1}$ = number of offspring of individual $i$ in generation $n$.
The key independence: the offspring counts $\{Z_i^{n+1}\}_{i\ge1}$ are
independent of $\mathcal{F}_n$ (they are ``new'' randomness in generation $n+1$).

\paragraph{Solution.}

\textbf{(i) Adaptedness.}
$X_n$ is determined by the first $n$ generations' offspring:
\begin{align*}
  X_1 &= Z_1^1 \quad\text{(from the single ancestor)}, \\
  X_2 &= \sum_{i=1}^{X_1}Z_i^2, \quad\ldots
\end{align*}
So $X_n\in\mathcal{F}_n$, and hence $M_n=\mu^{-n}X_n\in\mathcal{F}_n$.

\textbf{(ii) Integrability.}
We show by induction that $E[X_n]=\mu^n$, which gives $E[M_n]=1<\infty$.

\emph{Base case}: $E[X_0]=1=\mu^0$.

\emph{Inductive step}: Assume $E[X_n]=\mu^n$. Then
\[
  E[X_{n+1}] = E\!\left[\sum_{i=1}^{X_n}Z_i^{n+1}\right]
  = E\!\left[E\!\left[\sum_{i=1}^{X_n}Z_i^{n+1}\,\Bigg|\,X_n\right]\right].
\]
Given $X_n=k$, the sum $\sum_{i=1}^{k}Z_i^{n+1}$ consists of $k$ i.i.d.\ terms
each with mean $\mu$, so $E[\sum_{i=1}^{k}Z_i^{n+1}]=k\mu$. Hence
\[
  E\!\left[\sum_{i=1}^{X_n}Z_i^{n+1}\,\Bigg|\,X_n\right] = X_n\mu,
\]
and
\[
  E[X_{n+1}] = E[X_n\mu] = \mu E[X_n] = \mu\cdot\mu^n = \mu^{n+1}.
\]

\textbf{(iii) Martingale property.}
We compute $E[M_{n+1}\mid\mathcal{F}_n]$:
\begin{align*}
  E[M_{n+1}\mid\mathcal{F}_n]
  &= \mu^{-(n+1)}E[X_{n+1}\mid\mathcal{F}_n] \\
  &= \mu^{-(n+1)}E\!\left[\sum_{i=1}^{X_n}Z_i^{n+1}\,\Bigg|\,\mathcal{F}_n\right].
\end{align*}
Since $X_n$ is $\mathcal{F}_n$-measurable and $\{Z_i^{n+1}\}_{i\ge1}$ are
independent of $\mathcal{F}_n$, we can condition on $X_n=k$ and use:
\[
  E\!\left[\sum_{i=1}^{X_n}Z_i^{n+1}\,\Bigg|\,\mathcal{F}_n\right]
  = \sum_{k=0}^{\infty}\mathbf{1}_{\{X_n=k\}}\,E\!\left[\sum_{i=1}^{k}Z_i^{n+1}\right]
  = \sum_{k=0}^{\infty}\mathbf{1}_{\{X_n=k\}}\cdot k\mu
  = X_n\mu.
\]
(This is a standard application of: if $N$ is $\mathcal{F}_n$-measurable and
$\{W_i\}$ are i.i.d.\ independent of $\mathcal{F}_n$, then
$E[\sum_{i=1}^{N}W_i\mid\mathcal{F}_n]=N\cdot EW_1$.) Therefore,
\[
  E[M_{n+1}\mid\mathcal{F}_n]
  = \mu^{-(n+1)}\cdot X_n\mu
  = \mu^{-n}X_n = M_n.
\]
Hence $\{M_n\}$ is a martingale. $\square$

\paragraph{Remark.}
The martingale $\{M_n\}$ converges a.s.\ (by the martingale convergence theorem,
since $M_n\ge0$). The limit $M_\infty=\lim_{n\to\infty}\mu^{-n}X_n$ is related
to the extinction probability: if $\mu\le1$, then $M_\infty=0$ a.s.

% ======================================================================
\section*{Problem 5 (Superharmonic Functions)}

\paragraph{Problem statement.}
A function $f:\R^d\to\R$ is \emph{superharmonic} if, for all $x\in\R^d$ and all
$r>0$,
\[
  f(x) \ge \frac{1}{\vol(B(0,r))}\int_{B(x,r)}f(y)\,dy.
\]
Let $\{Z_n\}$ be i.i.d., uniformly distributed on $B(0,1)$. Set $S_0=0$ and
$S_n=S_{n-1}+Z_n$. Define $X_n=f(S_n)$. Show $\{X_n\}$ is a supermartingale
with respect to $\mathcal{F}_n=\sigma(Z_1,\ldots,Z_n)=\sigma(S_1,\ldots,S_n)$.

\paragraph{Solution.}

\textbf{(i) Adaptedness.}
$X_n=f(S_n)$. Since $S_n=Z_1+\cdots+Z_n$ is $\sigma(Z_1,\ldots,Z_n)=\mathcal{F}_n$-measurable,
so is $X_n=f(S_n)$.

\textbf{(ii) Integrability.}
We assume $f$ is integrable enough (e.g., $E|f(S_n)|<\infty$ for all $n$).
This holds, e.g., if $f$ is bounded.

\textbf{(iii) Supermartingale property.}
We compute $E[X_{n+1}\mid\mathcal{F}_n]$.

Since $S_{n+1}=S_n+Z_{n+1}$ and $Z_{n+1}$ is independent of $\mathcal{F}_n$
(it is the ``new'' increment), and $S_n$ is $\mathcal{F}_n$-measurable:
\[
  E[X_{n+1}\mid\mathcal{F}_n]
  = E[f(S_n+Z_{n+1})\mid\mathcal{F}_n].
\]
Condition on $\mathcal{F}_n$: since $Z_{n+1}\perp\mathcal{F}_n$ and $S_n$ is
$\mathcal{F}_n$-measurable, the conditional distribution of $S_n+Z_{n+1}$ given
$\mathcal{F}_n$ is the distribution of $S_n+Z_{n+1}$ with $S_n$ frozen at its
current value. That is, for fixed $s=S_n(\omega)$,
\[
  E[f(S_n+Z_{n+1})\mid\mathcal{F}_n](\omega)
  = E[f(s+Z_{n+1})]
  = \int_{B(0,1)} f(s+z)\,\frac{dz}{\vol(B(0,1))}.
\]
Now make the substitution $y=s+z$ (i.e., $z=y-s$). As $z$ ranges over
$B(0,1)=\{z:|z|\le1\}$, $y$ ranges over $\{y:|y-s|\le1\}=B(s,1)$. The Jacobian
is $1$. Hence:
\[
  \int_{B(0,1)}f(s+z)\,\frac{dz}{\vol(B(0,1))}
  = \frac{1}{\vol(B(s,1))}\int_{B(s,1)}f(y)\,dy
\]
where we use that $\vol(B(s,1))=\vol(B(0,1))$ (both are balls of radius 1 in $\R^d$,
just centered differently; their volume is the same).

By the \emph{superharmonic} property of $f$ applied at the point $x=s=S_n$ with
$r=1$:
\[
  f(s) \ge \frac{1}{\vol(B(s,1))}\int_{B(s,1)}f(y)\,dy.
\]
Therefore,
\[
  E[X_{n+1}\mid\mathcal{F}_n]
  = \frac{1}{\vol(B(S_n,1))}\int_{B(S_n,1)}f(y)\,dy
  \le f(S_n) = X_n.
\]
This holds a.s.\ (since it holds for every fixed realization of $S_n$).

Hence $E[X_{n+1}\mid\mathcal{F}_n]\le X_n$ a.s., so $\{X_n\}$ is a supermartingale.
$\square$

\paragraph{Remark.}
If $f$ is harmonic ($\ge$ replaced by $=$), then $\{X_n\}$ becomes a martingale.
Superharmonic functions correspond to supermartingales, and subharmonic to
submartingales --- this is the deep connection between potential theory and
martingale theory.

% ======================================================================
\section*{Problem 6}

\paragraph{Problem statement.}
Let $\{X_n\}$ be a DTMC on $\mathcal{X}$ with transition probabilities $\{p_{x,y}\}$.
Suppose $f:\mathcal{X}\to\R$ is bounded and satisfies
\[
  \sum_{y\in\mathcal{X}}p_{x,y}f(y)\le\lambda f(x)
\]
for some $\lambda>0$ and all $x\in\mathcal{X}$. Show that
$\{N_n\}:=\{\lambda^{-n}f(X_n)\}$ is a supermartingale with respect to
$\mathcal{F}_n=\sigma(X_0,\ldots,X_n)$.

\paragraph{Interpretation.}
The condition $\sum_y p_{x,y}f(y)\le\lambda f(x)$ says: the expected value of
$f$ after one step from $x$ is at most $\lambda f(x)$. This is the \emph{Foster-Lyapunov}
type condition. When we discount by $\lambda^{-n}$, the expected value after the
discount is non-increasing.

\paragraph{Solution.}

\textbf{(i) Adaptedness.}
$N_n=\lambda^{-n}f(X_n)$ is $\mathcal{F}_n$-measurable (it is a function of $X_n$,
which is $\mathcal{F}_n$-measurable).

\textbf{(ii) Integrability.}
Since $f$ is bounded, say $|f|\le M$, we have $|N_n|\le\lambda^{-n}M<\infty$, so
$N_n$ is bounded and integrable.

\textbf{(iii) Supermartingale property.}
\begin{align*}
  E[N_{n+1}\mid\mathcal{F}_n]
  &= \lambda^{-(n+1)}\,E[f(X_{n+1})\mid\mathcal{F}_n].
\end{align*}
By the Markov property:
\[
  E[f(X_{n+1})\mid\mathcal{F}_n]
  = E[f(X_{n+1})\mid X_n]
  = \sum_{y\in\mathcal{X}}p_{X_n,y}\,f(y).
\]
By the given condition applied at $x=X_n$:
\[
  \sum_{y\in\mathcal{X}}p_{X_n,y}\,f(y) \le \lambda\,f(X_n).
\]
Therefore,
\[
  E[N_{n+1}\mid\mathcal{F}_n]
  \le \lambda^{-(n+1)}\cdot\lambda\,f(X_n)
  = \lambda^{-n}f(X_n) = N_n.
\]
Hence $E[N_{n+1}\mid\mathcal{F}_n]\le N_n$ a.s., so $\{N_n\}$ is a supermartingale.
$\square$

\paragraph{Remark.}
If $\lambda=1$ and the inequality is an equality $\sum_y p_{x,y}f(y)=f(x)$, then
$f$ is a \emph{harmonic function} for the chain and $\{f(X_n)\}$ is a martingale.
If $\lambda=1$ and $\le$ holds, $f$ is superharmonic and $\{f(X_n)\}$ is a supermartingale.
Problem 6 generalizes this to arbitrary $\lambda>0$.

% ======================================================================
\section*{Problem 7 (P\'olya's Urn)}

\paragraph{Setup.}
Start with $r$ red and $b$ blue balls (total $r+b$, both $r,b>0$). At each step:
draw a ball uniformly at random, note its color, replace it, add a new ball of
the same color. $R_n$ = \#red balls after $n$ draws.

\subsection*{Part (a): $Y_n=R_n/(n+r+b)$ is a martingale}

\paragraph{Solution.}

\textbf{Setup.} After $n$ steps, total balls $= n+r+b$, red balls $= R_n$, blue balls $= n+r+b-R_n$.

\textbf{Adaptedness.} $Y_n=R_n/(n+r+b)$ is $\mathcal{F}_n=\sigma(R_1,\ldots,R_n)$-measurable.

\textbf{Integrability.} $0\le Y_n\le1$, so $E|Y_n|<\infty$.

\textbf{Martingale property.} At step $n+1$, given $\mathcal{F}_n$ (so $R_n$ is known):
\begin{itemize}
  \item With probability $\dfrac{R_n}{n+r+b}$, a red ball is drawn, so $R_{n+1}=R_n+1$.
  \item With probability $\dfrac{(n+r+b)-R_n}{n+r+b}$, a blue ball is drawn, so $R_{n+1}=R_n$.
\end{itemize}
Therefore,
\begin{align*}
  E[R_{n+1}\mid\mathcal{F}_n]
  &= (R_n+1)\cdot\frac{R_n}{n+r+b}
    + R_n\cdot\frac{(n+r+b-R_n)}{n+r+b} \\
  &= \frac{R_n^2+R_n + R_n(n+r+b)-R_n^2}{n+r+b} \\
  &= \frac{R_n(n+r+b) + R_n}{n+r+b} \\
  &= \frac{R_n(n+r+b+1)}{n+r+b}.
\end{align*}
Hence,
\[
  E[Y_{n+1}\mid\mathcal{F}_n]
  = \frac{E[R_{n+1}\mid\mathcal{F}_n]}{(n+1)+r+b}
  = \frac{R_n(n+r+b+1)}{(n+r+b)(n+r+b+1)}
  = \frac{R_n}{n+r+b}
  = Y_n.
\]
So $\{Y_n\}$ is a martingale. $\square$

\subsection*{Part (b): $E[(T+2)^{-1}]=1/4$ when $r=b=1$}

\paragraph{Setup.}
With $r=b=1$, the urn starts with 1 red and 1 blue ball (2 total). $T$ = number of
draws until the first blue ball.

\paragraph{Step 1: Distribution of $T$.}
At each draw, suppose the current number of red balls is $k$ and the total number of
balls is $k+1$ (since only red balls were added so far, having drawn all red for the
first $k-1$ draws). The probability of drawing blue (the only blue ball) is $\frac{1}{k+1}$.

More carefully:
\begin{itemize}
  \item At draw 1: urn has 1R, 1B. $P(\text{draw B})=\frac{1}{2}$, so $P(T=1)=\frac{1}{2}$.
  \item If draw 1 is red: urn becomes 2R, 1B (total 3). $P(\text{draw B})=\frac{1}{3}$.
  \item If draws 1,2 are red: urn becomes 3R, 1B (total 4). $P(\text{draw B at step 3})=\frac{1}{4}$.
\end{itemize}
In general, $T=k$ means the first $k-1$ draws were all red and draw $k$ is blue:
\[
  P(T=k)
  = \underbrace{\frac{1}{2}\cdot\frac{2}{3}\cdots\frac{k-1}{k}}_{\text{prob.\ first }k-1\text{ are red}}
    \cdot\underbrace{\frac{1}{k+1}}_{\text{draw }k\text{ is blue}}.
\]
The numerator of the first product is $(k-1)!$ and the denominator is
$2\cdot3\cdots k=\frac{k!}{1!}=k!$, so the first product is $\frac{(k-1)!}{k!}=\frac{1}{k}$.
Hence
\[
  \boxed{P(T=k) = \frac{1}{k(k+1)}}, \quad k=1,2,3,\ldots
\]
\emph{Verification}: $\sum_{k=1}^{\infty}\frac{1}{k(k+1)}
=\sum_{k=1}^{\infty}\!\left(\frac{1}{k}-\frac{1}{k+1}\right)=1$. $\checkmark$

\paragraph{Step 2: Computing $E[(T+2)^{-1}]$.}
\[
  E\!\left[\frac{1}{T+2}\right]
  = \sum_{k=1}^{\infty}\frac{1}{k+2}\cdot P(T=k)
  = \sum_{k=1}^{\infty}\frac{1}{k+2}\cdot\frac{1}{k(k+1)}
  = \sum_{k=1}^{\infty}\frac{1}{k(k+1)(k+2)}.
\]

\paragraph{Step 3: Partial fractions for $\frac{1}{k(k+1)(k+2)}$.}
We decompose:
\[
  \frac{1}{k(k+1)(k+2)} = \frac{A}{k}+\frac{B}{k+1}+\frac{C}{k+2}.
\]
Multiplying both sides by $k(k+1)(k+2)$:
\[
  1 = A(k+1)(k+2)+Bk(k+2)+Ck(k+1).
\]
\begin{itemize}
  \item $k=0$: $1=A\cdot1\cdot2$, so $A=\frac{1}{2}$.
  \item $k=-1$: $1=B(-1)(1)$, so $B=-1$.
  \item $k=-2$: $1=C(-2)(-1)$, so $C=\frac{1}{2}$.
\end{itemize}
Therefore $\dfrac{1}{k(k+1)(k+2)}=\dfrac{1/2}{k}-\dfrac{1}{k+1}+\dfrac{1/2}{k+2}$.

\paragraph{Step 4: Summing the series.}
\begin{align*}
  \sum_{k=1}^{\infty}\frac{1}{k(k+1)(k+2)}
  &= \frac{1}{2}\sum_{k=1}^{\infty}\frac{1}{k}
     - \sum_{k=1}^{\infty}\frac{1}{k+1}
     + \frac{1}{2}\sum_{k=1}^{\infty}\frac{1}{k+2}.
\end{align*}
Re-index: let $j=k+1$ in the second sum and $j=k+2$ in the third:
\begin{align*}
  &= \frac{1}{2}\sum_{k=1}^{\infty}\frac{1}{k}
     - \sum_{j=2}^{\infty}\frac{1}{j}
     + \frac{1}{2}\sum_{j=3}^{\infty}\frac{1}{j}.
\end{align*}
Split the first sum as $\sum_{k=1}^{\infty}\frac{1}{k}
=\frac{1}{1}+\frac{1}{2}+\sum_{j=3}^{\infty}\frac{1}{j}$, and the second sum as
$\sum_{j=2}^{\infty}\frac{1}{j}=\frac{1}{2}+\sum_{j=3}^{\infty}\frac{1}{j}$.
Let $S=\sum_{j=3}^{\infty}\frac{1}{j}$:
\begin{align*}
  &= \frac{1}{2}\!\left(1+\frac{1}{2}+S\right)
     - \!\left(\frac{1}{2}+S\right)
     + \frac{1}{2}S \\
  &= \frac{1}{2}+\frac{1}{4}+\frac{S}{2}
     - \frac{1}{2} - S
     + \frac{S}{2} \\
  &= \frac{1}{4} + \left(\frac{S}{2}-S+\frac{S}{2}\right)
   = \frac{1}{4} + 0
   = \frac{1}{4}.
\end{align*}
Hence $E\!\left[(T+2)^{-1}\right]=\dfrac{1}{4}$. $\square$

% ======================================================================
\section*{Problem 8 (Symmetric Random Walk: Gambler's Ruin)}

\paragraph{Setup.}
$Z_i$ i.i.d., $P(Z_i=1)=P(Z_i=-1)=\frac{1}{2}$. $X_0=0$, $X_n=\sum_{i=1}^nZ_i$.
Fix $a,b\in\mathbb{Z}_{>0}$ and $T=\min\{n\ge0:X_n\in\{-a,b\}\}$.

\paragraph{Preliminary: $T<\infty$ a.s.\ and $E[T]<\infty$.}
By standard results on the symmetric random walk on $\mathbb{Z}$, the walk
will hit any finite set of integers in finite expected time. In particular,
$T<\infty$ a.s.\ and $E[T]<\infty$ (in fact $E[T]=ab$ as we will show).

\subsection*{Part (a): $P(X_T=-a)=b/(a+b)$}

\paragraph{Martingale used.}
$\{X_n\}$ is itself a martingale, since $E[X_{n+1}\mid\mathcal{F}_n]
=X_n+E[Z_{n+1}]=X_n+0=X_n$.

\paragraph{Verifying OST conditions.}
For $n\le T$, we have $X_n\in(-a,b)$ (or $X_n\in\{-a,b\}$ when $n=T$), so
$|X_{n\wedge T}|\le\max(a,b)$. The stopped process $\{X_{n\wedge T}\}$ is bounded.
By the Bounded Optional Stopping Theorem (or Doob's OST), since the stopped process
is a bounded martingale and $T<\infty$ a.s.,
\[
  E[X_T] = E[X_0] = 0.
\]

\paragraph{Computing the probability.}
Let $p=P(X_T=b)$ so $P(X_T=-a)=1-p$. Since $X_T\in\{-a,b\}$:
\[
  0 = E[X_T] = b\cdot p + (-a)\cdot(1-p) = bp - a + ap = (a+b)p - a.
\]
Solving:
\[
  p = \frac{a}{a+b},
  \qquad
  P(X_T=-a) = 1-p = \frac{b}{a+b}.
  \qquad\square
\]

\subsection*{Part (b): $E[T]=ab$}

\paragraph{Key martingale.}
Consider $M_n=X_n^2-n$. We verify it is a martingale:
\begin{align*}
  E[M_{n+1}\mid\mathcal{F}_n]
  &= E[(X_n+Z_{n+1})^2-(n+1)\mid\mathcal{F}_n] \\
  &= E[X_n^2+2X_nZ_{n+1}+Z_{n+1}^2\mid\mathcal{F}_n] - (n+1) \\
  &= X_n^2 + 2X_n E[Z_{n+1}] + E[Z_{n+1}^2] - n - 1 \\
  &= X_n^2 + 0 + 1 - n - 1 = X_n^2-n = M_n.
\end{align*}
(We used $E[Z_{n+1}]=0$, $E[Z_{n+1}^2]=1$, and $Z_{n+1}\perp\mathcal{F}_n$.)

\paragraph{Applying OST.}
We apply OST to $\{M_n\}$ at time $T$. We need $E[M_T]$ to be well-defined.
The stopped martingale $\{M_{n\wedge T}\}$ satisfies:
\[
  |M_{n\wedge T}| = |X_{n\wedge T}^2-(n\wedge T)|
  \le \max(a,b)^2 + n\wedge T
  \le \max(a,b)^2 + T.
\]
Since $E[T]<\infty$ (established in Part (a) context and confirmed below), by
dominated convergence, $E[M_{n\wedge T}]\to E[M_T]$ as $n\to\infty$. Since
$\{M_{n\wedge T}\}$ is a martingale, $E[M_{n\wedge T}]=E[M_0]=0$ for all $n$.
Hence $E[M_T]=0$, i.e.,
\[
  E[X_T^2] = E[T].
\]

\paragraph{Computing $E[X_T^2]$.}
We use $P(X_T=b)=a/(a+b)$ and $P(X_T=-a)=b/(a+b)$:
\begin{align*}
  E[X_T^2]
  &= b^2\cdot\frac{a}{a+b} + a^2\cdot\frac{b}{a+b}
   = \frac{ab^2+a^2b}{a+b}
   = \frac{ab(a+b)}{a+b}
   = ab.
\end{align*}

\paragraph{Conclusion.}
\[
  E[T] = E[X_T^2] = ab.
  \qquad\square
\]

\paragraph{Remark on $E[T]<\infty$.}
To close the argument, we sketch why $E[T]<\infty$. Note $T\le T'$ where
$T'=\min\{n: X_n\notin(-a-b, a+b)\}$, and the random walk on a finite interval
has $E[T']<\infty$ by standard analysis (e.g., using the fact that at each step
there is a positive probability of exiting a box of size $a+b$ within $a+b$
steps). From $E[T]=ab$ itself being finite (once derived), we also get the
self-consistent conclusion.

\end{document}
